\small\begin{longtable}{p{0.13\textwidth}p{0.53\textwidth}p{0.23\textwidth}}
\caption{Referências selecionadas por década e acervo. Lexical não significa grafia mimética.}\\
\toprule Acervo / década & Documento e localização & Ocorrência / categoria\\\midrule\endfirsthead
\toprule Acervo / década & Documento e localização & Ocorrência / categoria\\\midrule\endhead
H04 / 1870 & A Imprensa Ytuana : jornal scientifico, litterario, noticiozo e industrial do Instituto do Novo-Mundo, ano 004, n. 183, 1879; 1879-09-13; PDF p. 2 & (Riso); rubrica \cite{r3doc1}\\
H04 / 1880 & A Imprensa Ytuana : jornal scientifico, litterario, noticiozo e industrial do Instituto do Novo-Mundo, ano 012, n. 101, 1887; 1887-05-13; PDF p. 2 & riso; lexical \cite{r3doc2}\\
H04 / 1890 & A Cidade de Ytu, ano 6, n. 383, 1898; 1898-07-03; PDF p. 1 & sorriso; lexical \cite{r3doc3}\\
H04 / 1900 & República, ano 009, n. 709, 1908; 1908-08-30; PDF p. 1 & Riso; lexical \cite{r3doc4}\\
H04 / 1910 & A Cidade de Ytu, ano 019, n. 1497, 1914; 1914-02-14; PDF p. 1 & riso; lexical \cite{r3doc5}\\
H04 / 1920 & A Cidade, ano 007, n. 606, 1929; 1929-02-10; PDF p. 1 & riso; lexical \cite{r3doc6}\\
H04 / 1930 & A Cidade, ano 010, n. 760, 1936; 1936-02-16; PDF p. 1 & riso; lexical \cite{r3doc7}\\
H04 / 1940 & O Bisturi, ano 016, n. 055, out. 1948; 1948; PDF p. 4 & sorriso; lexical \cite{r3doc8}\\
H04 / 1950 & O Trabalhador, ano 007, n. 110, 1955; 1955; PDF p. 1 & riso; lexical \cite{r3doc9}\\
H04 / 1960 & O Bisturi, ano 029, n. 112, nov. 1964; 1964-11; PDF p. 1 & risos; lexical \cite{r3doc10}\\
H04 / 1970 & O Bisturi, ano 041, n. 002, abr. 1976; 1976-04; PDF p. 5 & risos; lexical \cite{r3doc11}\\
H04 / 2000 & O Bisturi, ano 077, n. 005, jun. 2007; 2007-06; PDF p. 2 & sorriso; lexical \cite{r3doc12}\\
H04 / 2010 & O Bisturi : o jornal de estudantes de medicina da USP, ano 081, n. 001, jan./fev. 2011; 2011; PDF p. 1 & SORRISOS; lexical \cite{r3doc13}\\
H01 / 1920 & Correio da Semana; 1926; pasta 81, página lógica 39 & uma ironica gargalhada; lexical \cite{r3doc14}\\
H01 / 1930 & Correio da Semana; 1930; pasta 85, página lógica 33 & O Sorriso dos Velhos; lexical \cite{r3doc15}\\
H01 / 1960 & Correio da Semana; 1960; pasta 116, página lógica 161 & é de provocar o riso; lexical \cite{r3doc16}\\
H01 / 2000 & A Cidade de Jardinópolis; 2001; pasta 33, página lógica 209 & Ele cai na gargalhada; lexical \cite{r3doc17}\\
H01 / 2010 & A Cidade de Jardinópolis; 2016-01-09; pasta 49, página lógica 3; p. impressa 3; Só Rindo, Claudio Aguilhera & Hahaha...; mimética \cite{r3doc18}\\
H01 / 2020 & A Cidade de Jardinópolis; 2022-01-29; pasta 57, página lógica 21; p. impressa 5; Entrevista & kkkkk; mimética \cite{r3doc19}\\
H05 / 1840 & Diario Novo, n. 63; 1845-03-18; PDF p.1; Religião & os risos; lexical \cite{r3doc20}\\
H06 / 2020 & Correio Braziliense: Uma sintonia capital; 2025-04-07; PDF p.22; Diversão \& Arte; entrevista de Walter Lima Jr. por Ricardo Daehn & (risos); rubrica \cite{r3doc21}\\
H01 / 1970 & A Cidade de Jardinópolis, n.24; 1972-10-08; capa; pasta 4, página lógica 107; Divertindo-se às Margens do Rio Paraná & muitos risos, muitas alegrias; lexical \cite{r3doc22}\\
H01 / 1940 & Correio da Semana; 1940; pasta 95, página lógica 28 & reinado do riso; lexical \cite{r3doc23}\\
H01 / 1950 & Correio da Semana; 1950; pasta 106, página lógica 122 & se inspira lastima nem riso; lexical \cite{r3doc24}\\
H01 / 1980 & A Cidade de Jardinópolis; 1983; pasta 15, página lógica 49 & provocar o riso ao assistente; lexical \cite{r3doc25}\\
H01 / 1990 & A Cidade de Jardinópolis; 1990; pasta 22, página lógica 5 & romance e riso; lexical \cite{r3doc26}\\
H05 / 1920 & Diário Nacional, n. 290; 1928; 17/06/1928; p. 4; discurso de Mattos Pimenta & (risos, muito bem); rubrica \cite{r3doc27}\\
H05 / 1930 & O Tico-Tico, n. 1569; 1935; 30/10/1935; capa, último quadro; PDF p. 1 & Quá quá quá...; mimética \cite{r3doc28}\\
H05 / 1940 & Diário Carioca, n. 3902; 1941; 08/03/1941; p. 13; folhetim As memórias de Sherlock Holmes & Ah! ah!; mimética \cite{r3doc29}\\
\bottomrule\end{longtable}\normalsize
